% --- Plik: thesis-main.tex ---
% Druk dwustronny, rozmiar A4

\documentclass[11pt, a4paper, twoside]{report}

% Załadowanie ustawień i makr z osobnych plików
% --- Plik: thesis-packages.tex ---

%% Do kompilacji używaj XeLaTeX!
\usepackage{fontspec}

% Domyślna czcionka dokumentu (zamiennik dla Helvetica wg wytycznych)
\setmainfont{Arial}
\setsansfont{Arial}
\renewcommand{\familydefault}{\sfdefault}

% Definiowanie niestandardowych czcionek uczelnianych
% (Nazwy w nawiasach klamrowych muszą odpowiadać tym zainstalowanym w Twoim systemie)
\newfontfamily\adagioslab{Adagio_Slab-Regular}
\newfontfamily\adagioslablight{Adagio_Slab-Light}

\usepackage{polyglossia}
\setdefaultlanguage{polish}
\setotherlanguage{english}

% Ustawienia marginesów: wewnętrzny 30 mm, zewnętrzny 20 mm, górny i dolny 25 mm [3]
\usepackage[inner=30mm, outer=20mm, top=25mm, bottom=25mm]{geometry}

% Interlinia 1.15 [3]
\usepackage{setspace}
\setstretch{1.15}

% Pakiety dodatkowe
\usepackage{graphicx}
\usepackage{fancyhdr}
\usepackage{url}
\usepackage{xcolor}    % Pakiet do obsługi kolorów
\usepackage{hyperref}  % Pakiet do obsługi hiperłączy

\hypersetup{
	colorlinks=true,   % Włącza kolorowanie tekstu zamiast rysowania ramek
	linkcolor=black,   % Odnośniki wewnętrzne (np. spis treści) pozostają czarne do druku
	citecolor=black,   % Odnośniki do bibliografii pozostają czarne
	urlcolor=blue      % Adresy URL i e-mail będą wyraźnie niebieskie
	%urlcolor=blue!80!black
}


\usepackage{titlesec}
\usepackage{caption} % Do formatowania podpisów
\usepackage{booktabs}
\usepackage[table]{xcolor} % Włącza zaawansowane kolorowanie tabel

% Pakiety do obsługi zaawansowanej matematyki
\usepackage{amsmath}
\usepackage{amsfonts}


\usepackage{listings} % Do obsługi kodów źródłowych programów (np. Python, C++, MATLAB)
% Kolory do kodów Pythona w Jupyter Notebook
\definecolor{ju_keyword}{rgb}{0.0, 0.5, 0.0}
\definecolor{ju_string}{rgb}{0.73, 0.13, 0.13}
\definecolor{ju_comment}{rgb}{0.25, 0.50, 0.50}
\definecolor{ju_builtin}{rgb}{0.0, 0.30, 0.60}
\definecolor{ju_func}{rgb}{0.0, 0.0, 1.0}
\definecolor{ju_tlo}{rgb}{0.99, 0.99, 0.99}
\definecolor{ju_linia}{rgb}{0.85, 0.85, 0.85}
\lstset{
	language=Python,
	alsoletter={_},
	backgroundcolor=\color{ju_tlo},
	basicstyle=\ttfamily\small\color{black},
	keywordstyle=\bfseries\color{ju_keyword},
	commentstyle=\itshape\color{ju_comment},
	stringstyle=\color{ju_string},
	morekeywords=[2]{print,len,range,str,int,float,list,dict,set,tuple,zip,enumerate,abs,min,max,sum},
	keywordstyle=[2]\color{ju_builtin},
	emph={metoda_trapezow,metoda_simpsona,wykres_bledow,moja_funkcja},
	emphstyle=\color{ju_func},
	frame=single,
	rulecolor=\color{ju_linia},
	numbers=left,
	numberstyle=\tiny\color{gray!70},
	breaklines=true,
	columns=flexible,
	keepspaces=true,
	xleftmargin=15pt,
	xrightmargin=5pt,
	literate={ą}{{\k{a}}}1 {ć}{{\'c}}1 {ę}{{\k{e}}}1 {ł}{{\l}}1 {ń}{{\'n}}1 {ó}{{\'o}}1 {ś}{{\'s}}1 {ź}{{\'z}}1 {ż}{{\.z}}1
}

% Ciągła numeracja wzorów przez całą pracę
\counterwithout{equation}{chapter}

% Ciągła numeracja tabel przez całą pracę
\counterwithout{table}{chapter}

% Ciągła numeracja rysunków przez całą pracę
\counterwithout{figure}{chapter}

\AtBeginDocument{
    \counterwithout{lstlisting}{chapter} % Ciągła numeracja listingów przez całą pracę
	\renewcommand{\lstlistingname}{Kod źródłowy}   % Zmienia "Listing 1" na "Kod źródłowy 1"
}


% Konfiguracja stylu stron, nagłówków i stopek (Pakiet fancyhdr)
% Numeracja stron: na dole, po zewnętrznej stronie, druk dwustronny
\pagestyle{fancy}
\fancyhf{}
\fancyfoot[LE,RO]{\thepage}        % Bieżący numer strony: lewa parzysta (LE), prawa nieparzysta (RO)
\renewcommand{\headrulewidth}{0pt} % Całkowite usunięcie poziomej linii z nagłówka

% Zapewnienie spójności formatowania na stronach otwierających rozdziały 
% Domyślnie LaTeX wymusza na początku każdego rozdziału styl 'plain' (numer na środku).
% Poniższy blok nadpisuje styl 'plain', przenosząc numer na krawędzie boczne (LE/RO).
\fancypagestyle{plain}{%
	\fancyhf{}%
	\fancyfoot[LE,RO]{\thepage}%
	\renewcommand{\headrulewidth}{0pt}%
}

% Rozwiązanie konfliktu z pakietem formatowania tytułów (titlesec)
% Pakiet 'titlesec' często nadpisuje zachowanie komendy \chapter i ignoruje styl 'plain'.
% Poniższa instrukcja bezwzględnie wymusza stosowanie stylu 'fancy' na stronach z \chapter.
\usepackage{titlesec}
\assignpagestyle{\chapter}{fancy}


% Formatowanie nagłówków rozdziałów (pogrubiona bezszeryfowa: 14pt, 13pt, 12pt) [5]
\titleformat{\chapter}[hang]{\normalfont\sffamily\bfseries\fontsize{14}{16}\selectfont}{\thechapter.}{1em}{}
\titleformat{\section}{\normalfont\sffamily\bfseries\fontsize{13}{15}\selectfont}{\thesection}{1em}{}
\titleformat{\subsection}{\normalfont\sffamily\bfseries\fontsize{12}{14}\selectfont}{\thesubsection}{1em}{}

% Konfiguracja bibliografii (BibLaTeX + Biber)
\usepackage[
backend=biber,          % Silnik sortujący (Biber jest nowoczesny i obsługuje UTF-8)
style=numeric,          % Styl cytowania: numeryczny [1], [2] (alternatywa: alphabetic lub authoryear)
sorting=nty,            % Sortowanie w spisie: n-nazwisko, t-tytuł, y-rok
defernumbers=true       % Prawidłowe numerowanie przy podziale na sekcje
]{biblatex}

% Wskazanie pliku bazy danych zawierającego wpisy bibliograficzne
\addbibresource{thesis-references.bib}

% --- Plik: thesis-macros.tex ---

% Przykładowe makro do szybkiego pisania symboli ułamków
\newcommand{\R}{\mathbb{R}}
\newcommand{\C}{\mathbb{C}}
% Tutaj możesz dodawać własne skróty komend

%Jeśli wybierasz klasyczny wygląd z wcięciami, dodaj to do preambuły:
\setlength{\parindent}{0.5cm} % Wymusza wcięcie dokładnie 0,5 cm dla każdego akapitu
\setlength{\parskip}{0pt}     % Wyłącza dodatkowe odstępy pionowe między akapitami

%Jeśli wolisz blokowy układ tekstu, dodaj do preambuły ten kod:
%\setlength{\parindent}{0pt}   % Całkowicie wyłącza wcięcia akapitowe
%\setlength{\parskip}{4pt}     % Ustawia dokładnie 4 punkty pustej przestrzeni przed każdym nowym akapitem

% Konfiguracja podpisów (tytułów) tabel i rysunków
\captionsetup{
	font={sf,footnotesize}, % sf = czcionka bezszeryfowa, footnotesize = odpowiednik ok. 9 pkt
	justification=raggedright, % Wymusza justowanie do lewej krawędzi
	singlelinecheck=false,     % Zapobiega automatycznemu wyśrodkowaniu krótkich podpisów
	labelsep=period            % Wstawia kropkę po "Tabela 1" zamiast dwukropka
}



\begin{document}
	
	% Tymczasowe wyłączenie domyślnej numeracji dla stron początkowych (stylu 'plain')
	\makeatletter
	\let\origpsplain\ps@plain
	\let\ps@plain\ps@empty
	\makeatother
	\pagestyle{empty}
	
	% ---------------- STRONA TYTUŁOWA (Strona 1) ----------------
	\begin{titlepage}
		\noindent
		% Rejestracja bloku tekstu, aby zmierzyć szerokość napisu "Politechnika Warszawska"
		\newbox{\pwbox}
		\sbox{\pwbox}{{\adagioslab\fontsize{24}{28}\selectfont Politechnika Warszawska}}
		
		\begin{minipage}[t]{0.75\textwidth}
			\vspace{0pt} % Wyrównanie do góry
			% Politechnika Warszawska - wymóg: Adagio_Slab 24 pt [4]
			\usebox{\pwbox}\\[0.5cm]
			% Nazwa wydziału - wymóg: Adagio_Slab 12 pt, z rozstrzeleniem dopasowanym do szerokości napisu Politechnika Warszawska [4]
			\makebox[\wd\pwbox][s]{{\adagioslab\fontsize{12}{14}\selectfont W\hfill Y\hfill D\hfill Z\hfill I\hfill A\hfill Ł\hfill \ \hfill M\hfill E\hfill C\hfill H\hfill A\hfill T\hfill R\hfill O\hfill N\hfill I\hfill K\hfill I}}
		\end{minipage}%
		\begin{minipage}[t]{0.25\textwidth}
			\vspace{0pt}
			\begin{flushright}
				% Godło z prawej strony o średnicy 25 mm [4]
				\includegraphics[width=25mm]{logopw.png}
			\end{flushright}
		\end{minipage}
		
		\vspace{1.5cm}
		\begin{center}
			% Typ pracy - wymóg: Adagio_Slab Light 43 pt [5]
			{\adagioslablight\fontsize{43}{50}\selectfont Praca dyplomowa}\\[0.5cm]
			{\adagioslablight\fontsize{43}{50}\selectfont magisterska}\\[1.5cm]
			
			% Dalsze teksty używają domyślnej czcionki dokumentu (zamiennik dla Helvetica Light) [5, 7]
			{\fontsize{12}{14}\selectfont na kierunku Automatyka, Robotyka i Informatyka Przemysłowa}\\[1.5cm]
			
			% Tytuł pracy bez przedrostka - 14 pt (zastępstwo Helvetica) [5, 7]
			{\fontsize{14}{16}\selectfont Projekt i implementacja schedulera systemów operacyjnych o złożoności O(1) }\\[2cm]
			
			% Imię i Nazwisko - 21 pt (zastępstwo Helvetica) [6, 7]
			{\fontsize{21}{24}\selectfont Kajetan Kontek}\\[0.5cm]
			{\fontsize{12}{14}\selectfont numer albumu 327749}\\[1.5cm]
			
			{\fontsize{12}{14}\selectfont promotor}\\[0.5cm]
			{\fontsize{12}{14}\selectfont dr inż. Bartłomiej Fajdek}
			
			% Miękkie wypchnięcie tekstu na dół strony
			\vfill
			{\fontsize{12}{14}\selectfont WARSZAWA 2027}
		\end{center}
	\end{titlepage}
	
% Pusta strona po stronie tytułowej (Strona 2)
\clearpage\thispagestyle{empty}\mbox{}\clearpage

% ---------------- STRESZCZENIE PL (Strona 3) ----------------

\clearpage % Rozpoczynamy nową stronę
\thispagestyle{plain} % Ukrywa nagłówki, zostawia sam numer strony na dole (zgodnie z PW)

\begingroup
% Wymuszenie czcionki 12pt oraz interlinii 14pt (standardowa pojedyncza interlinia dla 12pt)
\fontsize{12pt}{14pt}\selectfont 

% Wycentrowany tytuł pracy po polsku
\begin{center}
	\textbf{Streszczenie} \\[0.5cm]
	\textbf{Modułowy szablon pracy dyplomowej w systemie \LaTeX{}}
\end{center}

\vspace{0.5cm}

\noindent Wytyczne Rektora Politechniki Warszawskiej precyzują kilka rygorystycznych wymogów odnośnie streszczenia pracy dyplomowej. Streszczenie stanowi odrębny punkt w układzie pracy i musi znajdować się na \textbf{osobnej stronie}. Obowiązuje tu również \textbf{ścisła objętość i formatowanie}, ponieważ każde streszczenie, niezależnie od języka, musi zajmować dokładnie jedną stronę, przy czym należy w nim zastosować pojedynczy odstęp (interlinię) oraz czcionkę o rozmiarze 12 pkt. 

Jeśli chodzi o \textbf{wymagane elementy}, to poza samym tekstem podsumowującym pracę, w każdym streszczeniu obowiązkowo musi pojawić się tytuł pracy oraz zestaw słów kluczowych. W praktyce akademickiej jako zestaw słów kluczowych przyjmuje się zazwyczaj od 3 do 5 pojęć lub krótkich fraz, które najdokładniej oddają tematykę i zakres przeprowadzonych w pracy badań. Umieszcza się je pod tekstem streszczenia.

Same \textbf{wersje językowe} zależą bezpośrednio od tego, w jakim języku napisana jest cała praca. Dla pracy w języku polskim wymagane jest osobne streszczenie po polsku oraz drugie streszczenie po angielsku. Dla pracy w języku angielskim wymagane jest streszczenie po angielsku oraz streszczenie po polsku (w tej kolejności). Z kolei dla pracy w innym języku obcym wymagane jest streszczenie w języku, w którym napisano pracę, a ponadto streszczenie po polsku oraz streszczenie po angielsku. \textbf{Opcjonalnie}, jeśli zamierzasz wystąpić o wydanie odpisu dyplomu w tłumaczeniu na jeszcze inny język obcy, powinieneś dołączyć dodatkowe streszczenie napisane w tym właśnie języku.

Warto zadbać o to, aby styl pisania był zwięzły i odzwierciedlał strukturę pracy: cel, metody, wyniki i wnioski.

\vspace{1.0cm} % Zwiększony odstęp pionowy 
\noindent \textbf{Słowa kluczowe:} praca dyplomowa, układ modułowy, \LaTeX{}

\endgroup
\clearpage % Wymuszenie końca strony - gwarancja, że streszczenie zajmie DOKŁADNIE 1 stronę


% Pusta strona po streszczeniu polskim (Strona 4)
\clearpage\thispagestyle{empty}\mbox{}\clearpage

% ---------------- STRESZCZENIE EN (Strona 5) ----------------
\selectlanguage{english}
% Identyczny mechanizm dla wersji obcojęzycznej
\thispagestyle{plain}
\begingroup
\fontsize{12pt}{14pt}\selectfont 

% Wycentrowany tytuł pracy po angielsku
\begin{center}
	\textbf{Abstract} \\[0.5cm]
	\textbf{Modular diploma thesis template in LaTeX system}
\end{center}

\vspace{0.5cm}

\noindent Here goes the English version of your abstract. It must follow the exact same formatting rules: 12pt font size, single line spacing, and it must fit entirely on this single page.

\vspace{1.0cm}

\noindent \textbf{Keywords:} diploma thesis, modular layout, \LaTeX{}
\endgroup
\clearpage % Koniec strony dla abstractu

% Powrót do języka polskiego
\selectlanguage{polish}


% Pusta strona po streszczeniu angielskim (Strona 6)
\clearpage\thispagestyle{empty}\mbox{}\clearpage

% ---------------- SPIS TREŚCI (Strona 7) ----------------
\tableofcontents

% Pusta strona po spisie treści (Strona 8)
\clearpage\thispagestyle{empty}\mbox{}\clearpage

% ---------------- TREŚĆ GŁÓWNA PRACY (Strona 9) ----------------
% Przywrócenie standardowej numeracji (styl 'plain') dla pierwszych stron rozdziałów
\makeatletter
\let\ps@plain\origpsplain
\makeatother
% Włączenie widocznej numeracji stron od tego miejsca [4]
\pagestyle{fancy} 

% Dołączanie kolejnych rozdziałów z osobnych plików
\chapter{Istniejące algorytmy planowania czasu procesora w nowoczesnych systemach}
\label{ch:wprowadzenie}

\subsection{System Linux i architektura EEVDF}
W nowoczesnych wersjach jądra systemu Linux, domyślnym algorytmem planowania (ang. \textit{scheduler}) dla procesów standardowych jest algorytm Earliest Eligible Virtual Deadline First (EEVDF). Zastąpił on wieloletni standard znany jako Completely Fair Scheduler (CFS), rozwiązując fundamentalny problem balansu między sprawiedliwością przydziału czasu a opóźnieniami w obsłudze zadań interaktywnych \cite{1}, \cite{38}. 

Działanie algorytmu EEVDF opiera się na ścisłych podstawach matematycznych \cite{15} i operuje na kilku kluczowych wartościach przypisywanych każdemu zadaniu:
\begin{itemize}
	\item \textbf{Żądany kwant czasu} (\textit{requested time slice / slice})
	\item \textbf{Wirtualny czas wykonania} (\textit{virtual runtime / vruntime})
	\item \textbf{Opóźnienie} (\textit{lag})
	\item \textbf{Priorytet} (\textit{priority / weight})
\end{itemize}

Kiedy nowe lub powracające ze stanu uśpienia zadanie trafia do kolejki procesów gotowych, otrzymuje bazowy kwant czasu (\textit{slice}). Wartość ta oparta jest na zmiennych konfiguracyjnych systemu (np. \texttt{sysctl\_sched\_base\_slice}) i jest dynamicznie skalowana na podstawie wagi procesu (wynikającej z jego priorytetu \textit{nice}) \cite{19}, \cite{26}. Następnie, gdy zadanie zaczyna być wykonywane przez procesor, jego wirtualny czas wykonania (\texttt{vruntime}) stale rośnie. Wartość ta odzwierciedla rzeczywisty czas spędzony na procesorze, znormalizowany przez wagę zadania. Zliczanie tego czasu opiera się na precyzyjnych licznikach sprzętowych procesora (np. w architekturze amd64), jednak sama aktualizacja wirtualnego czasu i struktur danych jest w całości realizowana programowo przez jądro systemu operacyjnego, zazwyczaj podczas obsługi przerwań zegarowych \cite{18}, \cite{19}.

Na bazie początkowo przypisanego kwantu czasu oraz \texttt{vruntime}, system wylicza wirtualny termin zakończenia (\textit{virtual deadline}). Reprezentuje on moment, do którego proces powinien w idealnych warunkach skonsumować swój przydzielony czas. Algorytm na bieżąco oblicza również opóźnienie (\textit{lag}), będące różnicą pomiędzy czasem, jaki zadanie powinno było otrzymać w idealnym, proporcjonalnym modelu podziału czasu (tzw. \textit{fluid-flow model}), a czasem, który faktycznie wykorzystało. Jeżeli wartość \textit{lag} jest mniejsza od zera, oznacza to, że proces otrzymał więcej czasu procesora niż pozostałe procesy rywalizujące o czas procesora. W takim scenariuszu proces traci kwalifikowalność (\textit{ineligible}) i jest kategorycznie wykluczony z możliwości przydziału sprzętowego do momentu, aż ogólny upływ czasu w systemie sprawi, że jego \textit{lag} ponownie wzrośnie do wartości nieujemnej \cite{1}, \cite{15}.

\subsection{Wybór procesu i struktury danych}
Działanie każdego planisty w nowoczesnych systemach operacyjnych opiera się na sprzętowym mechanizmie przerwań czasowych (ang. \textit{timer interrupts}). Ponieważ procesor nie może jednocześnie wykonywać własnego kodu schedulera oraz kodu aplikacji użytkownika, wbudowany kontroler przerwań z dużą częstotliwością wstrzymuje wykonywane zadanie, przekazując sterowanie do jądra systemu \cite{13}. Daje to algorytmowi możliwość przeliczenia wskaźników i, w razie potrzeby, wywłaszczenia procesu.

Scheduler EEVDF przechowuje informacje o aktywnych procesach w zrównoważonym drzewie czerwono-czarnym (ang. \textit{Augmented Red-Black Tree}). W przeciwieństwie do naiwnego podejścia, drzewo to nadal sortuje węzły na podstawie wirtualnego czasu wykonania (\texttt{vruntime}), co układa zadania od lewej do prawej. Kluczową innowacją jest jednak rozszerzenie (\textit{augmentation}) każdego węzła o dodatkowe pole \texttt{min\_deadline}, które przechowuje najkrótszy termin zakończenia z danego węzła i wszystkich jego potomków. Dzięki temu struktura działa jednocześnie jako drzewo wyszukiwania binarnego oraz jako kopiec (min-heap) dla wartości \textit{deadline}. Architektura ta pozwala systemowi w optymalnym czasie $O(\log n)$ odrzucić całe poddrzewa i natychmiast odnaleźć zadanie o najwcześniejszym \textit{virtual deadline}, które jednocześnie spełnia warunek kwalifikowalności ($lag \ge 0$) i to ono zostaje wybrane do działania na procesorze \cite{4}, \cite{19}.

\subsection{Mechanizm Run to Parity}
Algorytm EEVDF w swojej czystej, teoretycznej postaci zakłada wyidealizowany model maszyny obliczeniowej. W rzeczywistości fizyczne procesory opierają swoją wydajność na hierarchii pamięci podręcznej (Cache L1/L2/L3), która przechowuje dane powiązane z aktualnie wykonywanym procesem. Niezmodyfikowany mechanizm doboru zadań oparty wyłącznie na \textit{deadline'ach} prowadziłby do zjawiska tzw. \textit{over-scheduling} – zbyt częstych zmian aktualnie działającego zadania. Skutkowałoby to znacznymi stratami cykli procesora na ciągłą wymianę linii pamięci podręcznej (\textit{cache miss}) oraz narzutem związanym z przełączaniem kontekstu.

Aby rozwiązać ten problem, w systemie Linux wprowadzono mechanizm \textit{Run to Parity} \cite{26}. Zgodnie z tą regułą, proces, który już przejął kontrolę nad jednostką obliczeniową, z reguły nie zostaje wywłaszczony przez inne, standardowe zadania posiadające wcześniejszy \textit{virtual deadline}, dopóki jego własne opóźnienie (\textit{lag}) nie spadnie do zera. Dopiero gdy proces "wyrówna" swój dług wobec systemu (osiągając parytet), zdejmuje się z niego ochronę, pozwalając na przejęcie procesora przez kolejne kwalifikujące się zadanie. Mechanizm ten gwarantuje dłuższe, nieprzerwane odcinki pracy pojedynczego wątku, co maksymalizuje wykorzystanie pamięci podręcznej i zwiększa ogólną przepustowość sprzętu (throughput), nie naruszając jednocześnie długoterminowych matematycznych gwarancji sprawiedliwości algorytmu \cite{38}.
\chapter{Istniejące algorytmy planowania czasu procesora w nowoczesnych systemach}
\label{ch:istniejace}

\section{System Linux i architektura EEVDF}
W nowoczesnych wersjach jądra systemu Linux, domyślnym algorytmem planowania (ang. \textit{scheduler}) dla procesów standardowych jest algorytm Earliest Eligible Virtual Deadline First (EEVDF). Zastąpił on wieloletni standard znany jako Completely Fair Scheduler (CFS), rozwiązując fundamentalny problem balansu między sprawiedliwością przydziału czasu a opóźnieniami w obsłudze zadań interaktywnych \cite{linuxdocEEVDF, lkmlSchedV612}. 

Działanie algorytmu EEVDF opiera się na ścisłych podstawach matematycznych \cite{stoica1995eevdf} i operuje na kilku kluczowych wartościach przypisywanych każdemu zadaniu:
\begin{itemize}
	\item \textbf{Żądany kwant czasu} (\textit{requested time slice / slice})
	\item \textbf{Wirtualny czas wykonania} (\textit{virtual runtime / vruntime})
	\item \textbf{Opóźnienie} (\textit{lag})
	\item \textbf{Priorytet} (\textit{priority / weight})
\end{itemize}

Kiedy nowe lub powracające ze stanu uśpienia zadanie trafia do kolejki procesów gotowych, otrzymuje bazowy kwant czasu (\textit{slice}). Wartość ta oparta jest na zmiennych konfiguracyjnych systemu (np. \texttt{sysctl\_sched\_base\_slice}) i jest dynamicznie skalowana na podstawie wagi procesu (wynikającej z jego priorytetu \textit{nice}) \cite{linuxFairC, lkmlRunToParity}. Następnie, gdy zadanie zaczyna być wykonywane przez procesor, jego wirtualny czas wykonania (\texttt{vruntime}) stale rośnie. Wartość ta odzwierciedla rzeczywisty czas spędzony na procesorze, znormalizowany przez wagę zadania. Zliczanie tego czasu opiera się na precyzyjnych licznikach sprzętowych procesora (np. w architekturze amd64), jednak sama aktualizacja wirtualnego czasu i struktur danych jest w całości realizowana programowo przez jądro systemu operacyjnego, zazwyczaj podczas obsługi przerwań zegarowych \cite{linuxSchedH, linuxFairC}.

Na bazie początkowo przypisanego kwantu czasu oraz \texttt{vruntime}, system wylicza wirtualny termin zakończenia (\textit{virtual deadline}). Reprezentuje on moment, do którego proces powinien w idealnych warunkach skonsumować swój przydzielony czas. Algorytm na bieżąco oblicza również opóźnienie (\textit{lag}), będące różnicą pomiędzy czasem, jaki zadanie powinno było otrzymać w idealnym, proporcjonalnym modelu podziału czasu (tzw. \textit{fluid-flow model}), a czasem, który faktycznie wykorzystało. Jeżeli wartość \textit{lag} jest mniejsza od zera, oznacza to, że proces otrzymał więcej czasu procesora niż pozostałe procesy rywalizujące o czas procesora. W takim scenariuszu proces traci kwalifikowalność (\textit{ineligible}) i jest kategorycznie wykluczony z możliwości przydziału sprzętowego do momentu, aż ogólny upływ czasu w systemie sprawi, że jego \textit{lag} ponownie wzrośnie do wartości nieujemnej \cite{linuxdocEEVDF, stoica1995eevdf}.

\subsection{Wybór procesu i struktury danych}
Działanie każdego planisty w nowoczesnych systemach operacyjnych opiera się na sprzętowym mechanizmie przerwań czasowych (ang. \textit{timer interrupts}). Ponieważ procesor nie może jednocześnie wykonywać własnego kodu schedulera oraz kodu aplikacji użytkownika, wbudowany kontroler przerwań z dużą częstotliwością wstrzymuje wykonywane zadanie, przekazując sterowanie do jądra systemu \cite{wikiScheduling}. Daje to algorytmowi możliwość przeliczenia wskaźników i, w razie potrzeby, wywłaszczenia procesu.

Scheduler EEVDF przechowuje informacje o aktywnych procesach w zrównoważonym drzewie czerwono-czarnym (ang. \textit{Augmented Red-Black Tree}). W przeciwieństwie do naiwnego podejścia, drzewo to nadal sortuje węzły na podstawie wirtualnego czasu wykonania (\texttt{vruntime}), co układa zadania od lewej do prawej. Kluczową innowacją jest jednak rozszerzenie (\textit{augmentation}) każdego węzła o dodatkowe pole \texttt{min\_deadline}, które przechowuje najkrótszy termin zakończenia z danego węzła i wszystkich jego potomków. Dzięki temu struktura działa jednocześnie jako drzewo wyszukiwania binarnego oraz jako kopiec (min-heap) dla wartości \textit{deadline}. Architektura ta pozwala systemowi w optymalnym czasie $O(\log n)$ odrzucić całe poddrzewa i natychmiast odnaleźć zadanie o najwcześniejszym \textit{virtual deadline}, które jednocześnie spełnia warunek kwalifikowalności ($lag \ge 0$) i to ono zostaje wybrane do działania na procesorze \cite{kernelInternalsEEVDF, linuxFairC}.

\subsection{Mechanizm Run to Parity}
Algorytm EEVDF w swojej czystej, teoretycznej postaci zakłada wyidealizowany model maszyny obliczeniowej. W rzeczywistości fizyczne procesory opierają swoją wydajność na hierarchii pamięci podręcznej (Cache L1/L2/L3), która przechowuje dane powiązane z aktualnie wykonywanym procesem. Niezmodyfikowany mechanizm doboru zadań oparty wyłącznie na \textit{deadline'ach} prowadziłby do zjawiska tzw. \textit{over-scheduling} – zbyt częstych zmian aktualnie działającego zadania. Skutkowałoby to znacznymi stratami cykli procesora na ciągłą wymianę linii pamięci podręcznej (\textit{cache miss}) oraz narzutem związanym z przełączaniem kontekstu.

Aby rozwiązać ten problem, w systemie Linux wprowadzono mechanizm \textit{Run to Parity} \cite{lkmlRunToParity}. Zgodnie z tą regułą, proces, który już przejął kontrolę nad jednostką obliczeniową, z reguły nie zostaje wywłaszczony przez inne, standardowe zadania posiadające wcześniejszy \textit{virtual deadline}, dopóki jego własne opóźnienie (\textit{lag}) nie spadnie do zera. Dopiero gdy proces "wyrówna" swój dług wobec systemu (osiągając parytet), zdejmuje się z niego ochronę, pozwalając na przejęcie procesora przez kolejne kwalifikujące się zadanie. Mechanizm ten gwarantuje dłuższe, nieprzerwane odcinki pracy pojedynczego wątku, co maksymalizuje wykorzystanie pamięci podręcznej i zwiększa ogólną przepustowość sprzętu (throughput), nie naruszając jednocześnie długoterminowych matematycznych gwarancji sprawiedliwości algorytmu \cite{lkmlSchedV612}.






\section{Systemy z grupy Windows}
W nowoczesnych systemach z rodziny Windows za przydział czasu procesora odpowiada planista (ang. \textit{scheduler}) zaprojektowany pierwotnie na potrzeby architektury Microsoft Windows NT. W przeciwieństwie do algorytmów opartych na sprawiedliwym podziale czasu, rozwiązanie firmy Microsoft bazuje na wywłaszczeniowym planowaniu wielozadaniowym, w którym absolutnym fundamentem jest rygorystyczny priorytet \cite{loyolaProcThreadScheduling}. Co istotne, jednostką planowania w systemie Windows nie jest proces, lecz pojedynczy wątek. Proces pełni jedynie rolę kontenera na zasoby (takie jak pamięć czy kontekst zabezpieczeń), podczas gdy o dostęp do jednostki obliczeniowej rywalizują bezpośrednio wątki \cite{msLearnProcessesThreads}.

Architektura planisty definiuje 32 odrębne poziomy priorytetów (od 0 do 31), które podzielone są na trzy główne kategorie \cite{solomon2000threadScheduling}:
\begin{itemize}
	\item \textbf{Poziom 0 (Zero-Page Thread):} Zarezerwowany wyłącznie dla systemowego wątku bezczynności (\textit{idle thread}), który uruchamia się tylko przy braku innych zadań. Zajmuje się on m.in. asynchronicznym zerowaniem wolnych stron pamięci fizycznej.
	\item \textbf{Poziomy 1--15 (Priorytety dynamiczne):} Przydzielane standardowym wątkom aplikacji użytkownika oraz niekrytycznym wątkom jądra. Podlegają one ciągłym, dynamicznym algorytmom modyfikacji (wzmacnianiu i degradacji).
	\item \textbf{Poziomy 16--31 (Priorytety czasu rzeczywistego):} Zarezerwowane dla absolutnie krytycznych operacji (np. przetwarzanie multimediów przy minimalnych opóźnieniach, sterowniki sprzętowe). Wątki na tych poziomach mają sztywny priorytet i nie podlegają dynamicznym zmianom \cite{msLearnPriorityBoosts}.
\end{itemize}

Przy tworzeniu wątku w paśmie dynamicznym przydzielana jest mu tzw. wartość bazowa priorytetu (\textit{base priority}), obliczana na podstawie klasy priorytetu procesu oraz priorytetu wątku wewnątrz tej klasy. Aby zapewnić wysoką responsywność interaktywną i zapobiegać zbyt długiemu oczekiwaniu zadań (\textit{starvation}), planista stosuje mechanizm dynamicznego wzmacniania priorytetu (\textit{priority boosting}). W systemie istnieje sześć głównych rodzajów zdarzeń, które wyzwalają tymczasowe podbicie priorytetu wątku \cite{bsodTutorialsQuantumBoost}:
\begin{enumerate}
	\item Zakończenie asynchronicznej operacji wejścia/wyjścia (I/O completion).
	\item Wybudzenie wątku przez obiekty synchronizacji (np. semafory, zdarzenia).
	\item Zakończenie oczekiwania na zasoby warstwy wykonawczej (Executive resources).
	\item Działanie w ramach procesu pierwszoplanowego (Foreground process).
	\item Wybudzenie w celu obsługi zdarzeń interfejsu graficznego (GUI windowing activity).
	\item Interwencja usługi \textit{Balance Set Manager}, chroniąca przed głodzeniem poprzez wywindowanie priorytetu zablokowanego wątku do wartości 15 \cite{duffyBalanceSetManager}.
\end{enumerate}

Po otrzymaniu wzmocnienia, kwant czasu (ang. \textit{quantum}) wątku ulega wyzerowaniu. Kiedy wątek w pełni wykorzysta nowo przydzielony kwant, system stosuje mechanizm degradacji (\textit{priority decay}) – jego priorytet jest zmniejszany o 1. Proces ten powtarza się przy każdym kolejnym wyczerpaniu kwantu, aż wartość spadnie z powrotem do nominalnego priorytetu bazowego \cite{bsodTutorialsQuantumBoost}.

\subsection{Wybór procesu i struktury danych}
Kernel utrzymuje niezależne struktury danych dla każdego procesora logicznego, zgrupowane wewnątrz bloku \textit{Kernel Processor Control Block} (KPRCB) \cite{reactosThrdschdC}. Serce tej struktury stanowią kolejki gotowości (\textit{Dispatcher Ready Queues}), zaimplementowane jako tablica 32 dwukierunkowych list powiązanych – po jednej dla każdego poziomu priorytetu. Nowe lub wybudzone wątki dodawane są na koniec listy odpowiadającej ich priorytetowi.

W celu optymalizacji wyboru następnego zadania, dyspozytor unika czasochłonnego iterowania po listach. Utrzymuje on pojedynczą, 32-bitową liczbę, zwaną \textit{Ready Summary}. Każdy bit tej zmiennej powiązany jest z jedną kolejką. Jeśli dana kolejka posiada przynajmniej jeden wątek gotowy do natychmiastowego uruchomienia, odpowiadający jej bit ustawiany jest na wartość 1. Jeżeli kolejka jest pusta, bit przyjmuje wartość 0 \cite{solomon2000threadScheduling}.

Dzięki takiej architekturze, problem wyboru następnego wątku sprowadza się do odnalezienia najstarszego ustawionego bitu (ang. \textit{most significant set bit}) w zmiennej \textit{Ready Summary}. Jest to operacja matematyczna, której złożoność czasowa wynosi $O(1)$. W architekturach sprzętowych z rodziny x86/amd64 poszukiwanie to realizowane jest za pomocą pojedynczej instrukcji procesora – BSR (\textit{Bit Scan Reverse}). Pozwala to w bardzo krótkim czasie zidentyfikować najwyższą niepustą kolejkę, pobrać z niej pierwszy oczekujący wątek i zainicjować przełączenie kontekstu \cite{reactosKeX}.


\section{Systemy z rodziny BSD na przykładzie FreeBSD}
W nowoczesnych wersjach systemu FreeBSD (począwszy od wydania 5.2) domyślnym podsystemem szeregowania zadań jest planista ULE \cite{roberson2003ule}. Zastąpił on historyczny algorytm wywodzący się z systemu 4.4BSD, który ze względu na globalną kolejkę uruchomieniową oraz konieczność okresowego przeliczania priorytetów wszystkich zadań w systemie (złożoność $O(n)$) nie skalował się efektywnie w środowiskach wieloprocesorowych (SMP) \cite{roberson2003ule, mckusick2014freebsd}. Nazwa ULE nie jest akronimem -- pochodzi od trzech ostatnich liter angielskiego słowa \textit{schedule} (powstałego po usunięciu znaku podkreślenia z nazwy pliku źródłowego \texttt{sched\_ule.c}) \cite{roberson2003ule, freebsdSchedUleC}. Analogicznie jak w systemie Windows, podstawową jednostką podlegającą szeregowaniu we FreeBSD nie jest cały proces (\texttt{struct proc}), lecz pojedynczy wątek jądra (\texttt{struct thread}), posiadający własną strukturę metadanych planisty \texttt{struct td\_sched} \cite{freebsdSchedUleC}.

Architektura FreeBSD definiuje 256 dyskretnych poziomów priorytetów (od 0 do 255). Zgodnie z klasyczną konwencją systemów UNIX, niższa wartość liczbowa oznacza wyższy priorytet wykonywania \cite{mckusick2014freebsd}. Przestrzeń ta podzielona jest na cztery główne klasy \cite{mckusick2014freebsd}:
\begin{itemize}
	\item \textbf{Poziomy 0--47 (Wątki przerwań i jądra):} Najwyższe priorytety nadawane podsystemom jądra, gwarantujące bezwzględne prawo wywłaszczania procesów przestrzeni użytkownika.
	\item \textbf{Poziomy 48--79 (Wątki czasu rzeczywistego):} Zadania o krytycznym znaczeniu czasowym (\textit{real-time}), których priorytety są nadawane bezpośrednio przez aplikacje i nie podlegają kwantowaniu czasu.
	\item \textbf{Poziomy 120--223 (Wątki podziału czasu):} Pasmo przeznaczone dla standardowych procesów użytkownika (\textit{timesharing}), w którym priorytet wyliczany jest dynamicznie przez planistę.
	\item \textbf{Poziomy 224--255 (Wątki bezczynności):} Zarezerwowane dla zadań wykonywanych wyłącznie przy braku jakichkolwiek innych gotowych wątków w systemie (\textit{idle threads}).
\end{itemize}

W momencie utworzenia nowego procesu za pomocą wywołania \texttt{fork()}, wątek potomny dziedziczy początkowy poziom priorytetu oraz skumulowaną historię interaktywności bezpośrednio po swoim procesie macierzystym (np. powłoce systemowej) \cite{roberson2003ule}. Sprawia to, że nowo powołane zadania zazwyczaj rozpoczynają pracę z wysokim priorytetem charakterystycznym dla aplikacji interaktywnych. W trakcie dalszego działania wątku w paśmie podziału czasu (120--223), jego priorytet ulega ciągłym zmianom w zależności od zachowania zadania oraz przypisanego mu przez użytkownika parametru \textit{nice} (z przedziału od -20 do +20) \cite{roberson2003ule, mckusick2014freebsd}.

Kluczowym mechanizmem modyfikującym priorytet w czasie pracy jest heurystyka oceny interaktywności, realizowana przez funkcję \texttt{sched\_interact\_score()} \cite{freebsdSchedUleC}. Planista monitoruje w strukturze \texttt{td\_sched} skumulowany czas dobrowolnego uśpienia wątku ($T_{\text{sleep}}$) oraz czas jego faktycznej pracy na procesorze ($T_{\text{run}}$) w obrębie ruchomego okna historii (ograniczonego stałą \texttt{SCHED\_SLP\_RUN\_MAX}) \cite{roberson2003ule, freebsdSchedUleC}. Oznaczając przez $M$ połowę maksymalnej skali interaktywności, wskaźnik $S$ wyznaczany jest zależnie od relacji obu czasów. Gdy czas snu przewyższa czas pracy ($T_{\text{sleep}} > T_{\text{run}}$), wskaźnik przyjmuje postać \cite{roberson2003ule}:

$$S = \left\lfloor \frac{M \cdot T_{\text{run}}}{T_{\text{sleep}}} \right\rfloor$$

Natomiast gdy dominuje czas obliczeń ($T_{\text{run}} \ge T_{\text{sleep}}$), wartość obliczana jest ze wzoru \cite{roberson2003ule}:

$$S = M + \left\lfloor \frac{M \cdot (T_{\text{run}} - T_{\text{sleep}})}{T_{\text{run}}} \right\rfloor$$

Im niższa jest wyliczona wartość $S$, tym wyższą interaktywnością cechuje się zadanie. Jeżeli wskaźnik $S$ spadnie poniżej zdefiniowanego progu, wątek zostaje uznany za interaktywny, a jego priorytet jest wyznaczany proporcjonalnie do wyniku interaktywności (z uwzględnieniem modyfikatora \textit{nice}) \cite{roberson2003ule}. Jeśli jednak proces (np. kompilator uruchomiony z interaktywnej powłoki) zacznie intensywnie obciążać jednostkę obliczeniową bez dobrowolnego zasypiania, wartość $T_{\text{run}}$ szybko zdominuje $T_{\text{sleep}}$. Przekroczenie progu interaktywności spowoduje degradację zadania do kategorii wsadowej (\textit{batch}). Dla takich wątków priorytet wewnątrz pasma podziału czasu rośnie (czyli staje się mniej korzystny) wraz z liczbą wyczerpywanych taktów zegara statystycznego (\texttt{statclock}) oraz dodatnią wartością \textit{nice} \cite{roberson2003ule, mckusick2014freebsd}.

\subsection{Wybór procesu i struktury danych}
Aby wyeliminować rywalizację o blokady między rdzeniami i osiągnąć złożoność czasową $O(1)$, scheduler ULE utrzymuje osobny zestaw kolejek uruchomieniowych \texttt{struct tdq} dla każdego procesora logicznego (w tablicy \texttt{tdq\_cpu}) \cite{freebsdSchedUleC}. Wewnątrz każdego kontenera (zestawu) \texttt{tdq} znajdują się trzy wyspecjalizowane struktury \texttt{struct runq}:
\begin{enumerate}
	\item \texttt{tdq\_realtime} -- przechowuje wątki przerwań, zadania czasu rzeczywistego oraz wątki podziału czasu zaklasyfikowane jako interaktywne \cite{mckusick2014freebsd, freebsdSchedUleC}.
	\item \texttt{tdq\_timeshare} -- gromadzi nieinteraktywne (wsadowe) wątki podziału czasu. Działa ona jako kolejka kalendarzowa (\textit{calendar queue}), wykorzystująca dwa rotujące indeksy: zapisu (\texttt{tdq\_idx}) oraz odczytu (\texttt{tdq\_ridx}). Wstawiane wątki umieszczane są w odległości od \texttt{tdq\_idx} zależnej od ich priorytetu, natomiast pobieranie odbywa się od pozycji \texttt{tdq\_ridx}, która przesuwa się z upływem czasu. Gwarantuje to, że każde zadanie wsadowe zostanie uruchomione podczas pełnego obrotu wskaźnika, co całkowicie zapobiega zjawisku głodzenia (\textit{starvation}) bez potrzeby globalnego przeliczania priorytetów \cite{mckusick2014freebsd, freebsdSchedUleC}.
	\item \texttt{tdq\_idle} -- gromadzi wątki bezczynności, uruchamiane wyłącznie wtedy, gdy dwie powyższe struktury są całkowicie puste \cite{freebsdSchedUleC}.
\end{enumerate}

Na najniższym poziomie każda z trzech instancji \texttt{struct runq} agreguje 256 poziomów priorytetów w 64 dwukierunkowe listy ogonowe (\texttt{rq\_queues[64]}). Samo odnalezienie niepustej kolejki wewnątrz wybranej struktury realizowane jest w czasie stałym $O(1)$ przez funkcję \texttt{runq\_choose()}, która odczytuje 64-bitowy wektor stanu \texttt{rq\_status}. Pozycja pierwszego ustawionego bitu bezpośrednio wskazuje indeks najważniejszej niepustej listy, z której pobierany jest pierwszy oczekujący wątek \cite{mckusick2014freebsd, freebsdRunqH}. Jest to bardzo podobne rozwiązanie optymalizacyjne jak w systemach z rodziny Windows.
% --- Plik: chapter3_design.tex ---

\chapter{Projekt i realizacja systemu}
\label{ch:projekt} % To jest unikalny znacznik tego rozdziału

Niniejszy rozdział stanowi zasadniczą część merytoryczną pracy dyplomowej. Autor szczegółowo opisuje w nim strukturę projektową urządzeń, układów lub algorytmów oraz ich praktyczną realizację. 

Treść powinna w sposób ustrukturyzowany dokumentować proces realizacji projektu od założeń wstępnych do efektu końcowego, uwzględniając:
\begin{itemize}
	\item opracowanie koncepcji oraz uzasadnienie wyboru stosowanych metod, narzędzi, technologii lub elementów składowych,
	
	\item dokumentację projektową i inżynierską (np. schematy, modele matematyczne, diagramy blokowe, algorytmy),
	
	\item przebieg prac implementacyjnych, realizacyjnych lub symulacyjnych oraz integrację poszczególnych modułów i podsystemów.
\end{itemize}

\noindent \textbf{Stylistyka:} tekst należy redagować w bezosobowej formie czasu przeszłego (np. \textit{„zaprojektowano”, „zaimplementowano”, „przeprowadzono dobór”}), zachowując charakter sprawozdania technicznego.


 % --- Plik: chapter4_experiments.tex ---
 
 \chapter{Wyniki badań}
 \label{ch:wyniki} % To jest unikalny znacznik tego rozdziału
 
 Niniejszy rozdział służy prezentacji, eksperymentalnej walidacji oraz krytycznej analizie uzyskanych wyników badawczych, pomiarowych lub symulacyjnych. 
 
Treść powinna w sposób obiektywny dokumentować działanie i charakterystykę opracowanego rozwiązania, uwzględniając:
 \begin{itemize}
 	\item opis metodyki badawczej, warunków przeprowadzenia eksperymentu lub scenariuszy testowych,
 	
 	\item opracowane i czytelnie sformatowane dane wynikowe zaprezentowane w formie tabel zbiorczych oraz wykresów,
 	
 	\item interpretację uzyskanych zależności, analizę błędów pomiarowych oraz odniesienie wyników do założeń początkowych i stanu wiedzy.
 \end{itemize}
 \noindent \textbf{Stylistyka:} w tym rozdziale nie należy używać określeń subiektywnych (np. \textit{„wyniki są bardzo dobre”, „system działa świetnie”}). Należy posługiwać się wyłącznie językiem faktów, liczb i konkretnych wskaźników inżynierskich, stosując bezosobową formę czasu przeszłego.


% --- Plik: chapter5_conclusions.tex ---

\chapter{Podsumowanie}
\label{ch:podsumowanie} % To jest unikalny znacznik tego rozdziału

Niniejszy rozdział stanowi formalne zamknięcie merytorycznej części pracy dyplomowej. Jego celem jest zwięzła synteza uzyskanych rezultatów oraz globalna ocena stopnia realizacji zadania inżynierskiego.

W rozdziale tym należy dokonać całościowej oceny zrealizowanych prac inżynierskich, opierając strukturę wniosków na następujących kluczowych aspektach:
\begin{itemize}
	\item jednoznacznej weryfikacji stopnia realizacji celów oraz tez sformułowanych we wstępie pracy,
	
	\item zwięzłym omówieniu najważniejszych osiągnięć o charakterze konstrukcyjnym, programistycznym bądź eksperymentalnym,
	
	\item obiektywnej ocenie ograniczeń wypracowanego rozwiązania oraz wskazaniu potencjalnych kierunków jego dalszej optymalizacji lub rozbudowy.
\end{itemize}
\noindent \textbf{Stylistyka:} wnioski powinny być sformułowane w sposób konkretny i zwięzły. Należy unikać dosłownego powtarzania zdań z poprzednich rozdziałów, skupiając się na ogólnej konkluzji z całej pracy.

\vspace{0.5cm}

\noindent \textbf{Więcej informacji na temat systemu \LaTeX{}:}
\begin{itemize}
	\item \url{https://www.overleaf.com/learn} -- bardzo przystępny, oficjalny tutorial i baza wiedzy na stronie platformy Overleaf,
	
	\item \url{https://www.latex-project.org/} -- oficjalna strona domowa projektu \LaTeX{}, zawierająca aktualności i dokumentację standardów,
	
	\item \url{https://www.tug.org/begin.html} -- prowadzony przez \textit{\TeX{} Users Group} obszerny zbiór odnośników do materiałów dla osób początkujących,
	
	\item \url{https://www.texstudio.org/} -- oficjalna strona domowa edytora TeXstudio, z której należy pobierać najnowsze, stabilne wersje programu,
	
	\item \url{https://www.xm1math.net/texmaker/} -- oficjalna strona domowa edytora Texmaker, z której można pobrać program oraz uzyskać dostęp do jego podręcznika użytkownika,

	\item \url{https://texstudio-org.github.io/} -- oficjalny, szczegółowy podręcznik użytkownika (\textit{User Manual}) wyjaśniający zaawansowane funkcje edytora TeXstudio,
	
	\item \url{https://tex.stackexchange.com/} -- największe na świecie forum społecznościowe typu Q\&A (pytania i odpowiedzi), na którym można znaleźć rozwiązanie niemal każdego problemu technicznego i błędu kompilacji.
\end{itemize}


\newpage

% ---------------- BIBLIOGRAFIA ----------------

% Zmiana nazwy nagłówka z domyślnego "Bibliografia" na "Literatura" 
\addcontentsline{toc}{chapter}{Literatura}
\printbibliography[title={Literatura}]

% Pusta strona po stronie Literatura
\clearpage\thispagestyle{empty}\mbox{}\clearpage

% ---------------- WYKAZ SYMBOLI I SKRÓTÓW ----------------

% Wykaz symboli i skrótów - jak nie ma, to zakomentować lub usunąć
\chapter*{Wykaz symboli i skrótów}
\addcontentsline{toc}{chapter}{Wykaz symboli i skrótów} % Opcjonalnie: dodaje wykaz do spisu treści

\begin{tabular}{cl} % 'cl' centruje symbol i dosuwa opis do lewej
	$\omega$ & prędkość kątowa [rad/s] \\
	$M$      & moment obrotowy [Nm] \\
	$U$      & napięcie zasilania [V] \\
	$I$      & natężenie prądu [A] \\
	$\theta$ & kąt obrotu wału silnika [$^\circ$ lub rad] \\
	$F$      & siła ciągu/nacisku [N] \\
	$k_p, k_i, k_d$ & nastawy (wzmocnienia) regulatora PID \\
\end{tabular}

\vspace{1cm} % Odstęp między tabelami

\begin{tabular}{ll} % 'll' wyrównuje obie kolumny do lewej, co przy skrótach wygląda lepiej
	ADC  & (ang. \textit{Analog-to-Digital Converter}) przetwornik analogowo-cyfrowy \\
	DOF  & (ang. \textit{Degrees of Freedom}) stopnie swobody \\
	IoT  & (ang. \textit{Internet of Things}) internet rzeczy \\
	PLC  & (ang. \textit{Programmable Logic Controller}) programowalny sterownik logiczny \\
\end{tabular}


% Pusta strona po stronie Wykaz symboli i skrótów
\clearpage\thispagestyle{empty}\mbox{}\clearpage

% ---------------- SPIS RYSUNKÓW ----------------

\listoffigures   

% Pusta strona po stronie Spis rysunków
\clearpage\thispagestyle{empty}\mbox{}\clearpage


% ---------------- SPIS TABEL ----------------

\listoftables  

% Pusta strona po stronie Spis tabel
\clearpage\thispagestyle{empty}\mbox{}\clearpage


% ---------------- SPIS ZAŁĄCZNIKÓW ----------------

% Spis załączników - jak nie ma załączników, to zakomentować lub usunąć
\chapter*{Spis załączników}
\addcontentsline{toc}{chapter}{Spis załączników} % Opcjonalnie: dodaje wykaz do spisu treści

\begin{enumerate}
	\item Analiza błędów -- pełny skrypt źródłowy zawierający zaimplementowane funkcje testowe oraz procedurę numerycznej analizy błędów metody trapezów.
	
	\item Specyfikacja techniczna stacji roboczej wykorzystanej w pracy.
	
	\item Dodatkowe zdjęcia stanowiska badawczego.
	
	%\item Dokumentacja katalogowa -- noty aplikacyjne producenta dotyczące wykorzystanego czujnika momentu obrotowego.
\end{enumerate}

% Pusta strona po stronie Literatura
\clearpage\thispagestyle{empty}\mbox{}\clearpage



% --- ROZPOCZĘCIE SEKCJI ZAŁĄCZNIKÓW ---
\appendix
\setcounter{chapter}{0} % <--- Wymusza start numeracji od zera
% Zmiana słowa "Dodatek" na "Załącznik"
\renewcommand{\appendixname}{Załącznik}

% Zmiana formatu numeracji z liter (A, B, C) na cyfry (1, 2, 3)
\renewcommand{\thechapter}{\arabic{chapter}}


\chapter*{Załącznik 1\\Analiza błędów}
\refstepcounter{chapter}

% To jest unikalny znacznik tego załącznika; w tekście pracy można odwołać się do konkretnego załącznika
\label{zal:analiza_bledow} 

% Tutaj wpisujemy treść pierwszego załącznika lub wklejamy kod przez \lstinputlisting

Zarządzenie Rektora PW~\cite{rektorPW2022} ściśle definiuje układ dokumentu, wskazując, że przedostatnim jego elementem jest „Spis załączników”, w którym  wymienia się te materiały, które są fizycznie dołączone na samym końcu dokumentu pracy dyplomowej.

Poniżej znajduje się pełny skrypt zawierający funkcje testowe oraz procedurę analizy błędów.

% Pusta strona po stronie Literatura
\clearpage\thispagestyle{empty}\mbox{}\clearpage


\chapter*{Załącznik 2\\Specyfikacja techniczna stacji roboczej wykorzystanej w pracy}
\refstepcounter{chapter}
\label{zal:specyfikacja} % To jest unikalny znacznik tego załącznika

\begin{table}[htbp]
	\centering 
	\label{tab:specyfikacja_komputera}
	\begin{tabular}{|l|p{9.5cm}|} 
		\hline % Pozioma linia na samej górze
		\textbf{Komponent} & \textbf{Parametry techniczne / Model} \\
		\hline % Pozioma linia pod nagłówkiem
		Procesor (CPU) & Intel Core i7-14700K (20 rdzeni, 28 wątków, taktowanie do 5.6 GHz) \\
		\hline
		Pamięć RAM & 64 GB DDR5 Kingston Fury (2 $\times$ 32 GB, 6000 MHz, CL30) \\
		\hline
		Karta graficzna (GPU) & NVIDIA GeForce RTX 4080 Super (16 GB GDDR6X, 10240 rdzeni CUDA) \\
		\hline
		Dysk twardy & SSD Samsung 990 Pro 2 TB NVMe M.2 PCIe Gen4 (odczyt do 7450 MB/s) \\
		\hline
		Płyta główna & MSI MAG Z790 TOMAHAWK WIFI \\
		\hline
		Zasilacz & Corsair RM850x 850W (Certyfikat 80 Plus Gold) \\
		\hline
		System operacyjny & Linux Ubuntu 24.04 LTS / Windows 11 Pro (64-bit) \\
		\hline
		Środowisko i biblioteki & Python 3.12, CUDA Toolkit 12.4, PyTorch 2.3, MATLAB R2024a \\
		\hline % Pozioma linia na samym dole
	\end{tabular}
\end{table}

% Pusta strona po stronie Literatura
\clearpage\thispagestyle{empty}\mbox{}\clearpage


\chapter*{Załącznik 3\\Dodatkowe zdjęcia stanowiska badawczego}
\refstepcounter{chapter}
\label{zal:zdjecia} % To jest unikalny znacznik tego załącznika

W niniejszym załączniku zgromadzono dodatkowe materiały graficzne, dokumentację fotograficzną stanowiska badawczego oraz szczegółowe wykresy przebiegów, które ze względu na swoje znaczne rozmiary i pomocniczy charakter nie zostały włączone bezpośrednio do zasadniczej części pracy. Zamieszczone poniżej materiały stanowią bezpośrednie uzupełnienie opisów eksperymentów zawartych w rozdziale~\ref{ch:wyniki} i umożliwiają pełniejszą, szczegółową weryfikację poprawności działania zaimplementowanego systemu.


\end{document}